\section{From an utterance to a grammaticality profile}
\label{sec:formalism}

The preceding section supplied diagnostic labels. This section states the model
they belong to without asking the reader to learn its mathematical
representation. The full specification, derivations, simulations, and notation
are in the formal supplement. The main idea is simple: a grammaticality claim
compresses several answers about a form--value pairing in a particular
community and situation. The answers can come apart, and their pattern matters
more than any single score.

The model separates a \term{state theory} from a \term{dynamics theory}. The
state theory asks what has to be true for a candidate to count as grammatical
now. The dynamics theory asks how learning, choice, competition, repair, and
forgetting could make that state persist or change. This separation blocks a
common shortcut: a story about how a pattern arose doesn't by itself establish
the pattern's present status, and a stable present pattern doesn't identify the
mechanism that produced it.

\subsection{Fix the situation before diagnosing the form}
\label{sec:conditioning}
\label{sec:identification-c}

Every diagnosis is relative to a \term{conditioning state}: the communicative
situation as the participants construe it, together with the community, variety,
register, or norm centre to which they orient. This state can include the
activity, medium, footing, attributed speaker identity, and relevant discourse
community. It isn't assumed to be obvious or externally given. Interlocutors can
construe the same encounter differently, and analysts can draw the boundary too
broadly or too narrowly.

For that reason, the conditioning state has to be specified independently of
the judgment the model is meant to explain. A study can't label every accepting
speaker as belonging to one variety and every rejecting speaker as belonging to
another after seeing the responses. It needs prior evidence from community
membership, usage, interactional cues, self-identification, or a registered
coding protocol. If an independently motivated partition turns an apparently
divided population into internally stable groups, the original conditioning
state was underspecified. If no such partition does, the disagreement is a
genuine population fact.

\subsection{Four constitutive questions}
\label{sec:objects}

OVMG treats constructions in the broad Construction Grammar sense: learned,
typed associations among form, structure, content, context, and conventional
constraints, at any grain from an inflectional cell to a clause pattern or an
intonational contour \autocite{fillmore1988, goldberg1995constructions,
sag2012sign}. A candidate utterance can normally be assembled in more than one
way from such resources. The model asks whether \emph{some} analysis
of the intended form and value passes four tests.

\begin{enumerate}
\item \textbf{Coverage: can the language build it?} There has to be at least one
  complete analysis of the form with the intended value. Word salad may fail
  here. A merely novel sentence need not: productive constructions can cover a
  token that has never previously occurred.

\item \textbf{Compatibility: do its operator instructions fit together?} The
  contributions that tell interlocutors how to update the conversation have to be
  jointly satisfiable. In \ungram{\mention{I've finished it yesterday}}, the
  present perfect and the deictic past-time adverb provide incompatible temporal
  instructions. More exposure to the same analysis can't average that conflict
  away.

\item \textbf{Saturation: are the contrasts required in this niche supplied?}
  A community can come to treat overt marking of a distinction as obligatory.
  On this account, obligatoriness isn't simply listed in advance. It's the
  historical result of situations in which an unmarked alternative repeatedly
  loses to an overtly marked one. A moribund distinction can later cease to be
  obligatory.

\item \textbf{Licensing: are the constructions used in the analysis established
  resources here?} A form can be analyzable, interpretable, and internally
  compatible without being part of the relevant repertoire. English
  \ungram{\mention{I have 25 years}} illustrates a head-and-complement
  convention of this kind. The intended age meaning is recoverable, but English
  licenses a different frame. No operator has to be mis-set for a construction
  to fail this test.
\end{enumerate}

These tests are ordered for diagnosis but aren't four degrees on one scale.
Coverage and compatibility are properties of analyses in a specified
inventory. Saturation records a path-dependent community convention. Licensing
is a population pattern over constructional resources. A candidate is
grammatical when at least one analysis passes all four; alternative failed
analyses don't contaminate a successful one. This scope matters for ambiguity,
coercion, and productive novelty.

The distinction between operator value and exponent choice cuts across the four
tests. \ungram{\mention{Goed}} expresses the intended past value but selects an
unlicensed exponent. Turkish suffix-harmony errors can have the same shape: the
plural or past value remains identifiable, while the chosen allomorph isn't
licensed in that phonological environment. Conversely, \ungram{\mention{depend
of}} is a head-specific complement convention, not an operator error. Such
cases show why neither closed-class status nor categorical rejection defines
the operator stratum.

\subsection{Population status and a speaker's evidence}
\label{sec:posterior-existence}
\label{sec:objective-G}

The four questions concern a population profile, but no speaker or analyst sees
that profile directly. Each observes a history of tokens, alternatives,
omissions, repairs, and contextual cues and forms an estimate from it. OVMG
keeps three things distinct: the community's current licensing
pattern, an individual speaker's evidence-based estimate of that pattern, and an
analyst's estimate from a possibly different sample.

This distinction adds an important second dimension to the familiar gradient
mean. Two forms can receive the same middling rating for quite different
reasons. One speaker may have little evidence either way; another community may
be sharply divided between stable users and stable non-users. The first is
epistemic uncertainty and should be sensitive to framing and new exposure. The
second is population heterogeneity and should instead appear as stable
between-speaker disagreement. Mean ratings alone erase the difference.

The model also distinguishes uncertainty about the constructional inventory
from uncertainty about licensing. If analysts don't know whether a productive
analysis exists, that's an uncertainty about coverage. Once a viable analysis
is fixed, sparse evidence about whether a community uses its nodes is a
licensing question. Treating every unfamiliar token as absent from the inventory
would wrongly predict that first-heard but readily productive sentences
begin as ungrammatical.

\subsection{What the profile is expected to predict}
\label{sec:projectible-use}
\label{sec:projectible-G}

The profile earns its theoretical role only if observing part of it licenses
expectations about cases not already used to define it. In plain terms, the
analysis should tell us more than which label to attach to yesterday's rating.
It should help predict how a form will be repaired, learned, transmitted,
softened by exposure, or destabilized over time.

The central projection contrasts weak evidence with strong contrary evidence.
A construction rare because its niche seldom occurs should remain uncertain
and comparatively movable under informative exposure. A form
absent despite many genuine opportunities, and repeatedly displaced by an
established competitor, should produce confident and framing-resistant
rejection. A processing-heavy but licensed sentence should improve with
practice without a corresponding change in production or repair evidence. A
moribund contrast should first lose evidential support and only later lose its
categorical treatment. These targets go beyond the tests used to assign the
initial profile and can fail independently.

Those failures have to change the analysis. The situation, niche, and target
projection should be fixed before the outcome is inspected. If a profile then
fails to predict held-out repair, exposure, transmission, or change, the
projection has to be narrowed, split, or abandoned. It can't be rescued by
redrawing the speech-community boundary after the responses are known.

The supporting explanation is deliberately weaker than a claim of
self-regulation. A stable profile might be the residue of history, the result of
ongoing transmission, or the effect of an active corrective process. OVMG calls
these stabilization, maintenance, and control, respectively. Evidence for one
doesn't establish the next. In particular, repair is only a candidate
controller until intervention shows that a deviation is registered and that the
response helps restore the threatened higher-level relation.

\subsection{Status and the feeling of status}
\label{sec:decision-readout}
\label{sec:feeling-new}

Speakers don't inspect a population state. They experience an utterance and
respond to it. OVMG represents that response with two ordinary-language
dimensions: how anomalous the token feels and how confident the speaker is in
that reaction. Low expected licensing can contribute to anomaly, but so can
processing load, lexical implausibility, garden paths, prescriptive pressure,
and misidentification of the speaker or situation.

Confidence supplies information that anomaly alone misses. A preempted gap may
elicit a confident \enquote{that's impossible}. A winnerless cell may elicit an
equally low rating but the phenomenology \enquote{I don't know how to say it}.
An independent-relative \mention{whose} may feel unfamiliar while remaining
movable under framing. Centre embedding can feel very bad even when the speaker
accepts the construction once its parse is made clear. These aren't verbal
glosses added after the fact: the model predicts different patterns of
test--retest stability, confidence, exposure response, production, and repair.

A community or institution may impose a practical acceptance boundary on this
graded evidence. That boundary belongs to a decision layer: changing the costs
of false acceptance and false rejection should change how a hearer responds
without changing the underlying licensing profile. A sharp switch in overt
response would test the decision process, not prove that grammatical
knowledge itself contains a numerical cutoff.

\subsection{Misalignment and the scope of a successful analysis}
\label{sec:scope-existential}
\label{sec:misalignment-conditioning}

Hearer and speaker can disagree about which situation is in force. A form
licensed in one dialect, genre, age group, or bilingual practice may be heard
through another community's expectations. The resulting anomaly can reflect a
real mismatch between two stable systems rather than uncertainty within either
one. Attribution matters too: a hearer who treats a token as belonging to
another community, a learner error, noise, or personal overload need not respond
as if it threatened the local convention.

Within a fixed situation, the existence of one successful analysis is enough.
This prevents the model from rejecting ordinary ambiguity merely because one
available parse fails. But the value has to be fixed independently enough to
make the claim vulnerable. If analysts are free to invent a new value whenever
an analysis would otherwise fail, coverage becomes vacuous. The intended value
should be constrained by context, paraphrase, comprehension, and the speaker's
subsequent commitments.

\subsection{Measurement requires converging channels}
\label{sec:measurement-C}

Status isn't directly observed, so no single observation defines it. The
empirical programme combines four channels:

\begin{itemize}
\item \textbf{Production} shows which available alternatives speakers choose
  in independently specified niches. A low corpus rate can reflect
  unavailability, low preference, or avoidance, so production alone can't
  identify licensing.
\item \textbf{Repair} shows how interlocutors respond to a putative violation.
  Understanding repair, update-oriented repair, overt correction, and
  prescriptive sanction have to be coded separately.
\item \textbf{Judgment and confidence} measure the subjective read-out rather
  than status directly. Repetition, framing, and test--retest designs show how
  stable that read-out is.
\item \textbf{Situation cues} support the analyst's and hearer's identification
  of the relevant community and activity. They can't be inferred solely from
  the target rating.
\end{itemize}

A strong test requires at least one non-judgment indicator of licensing, an
independent measure of whether changing an operator value changes the public
update, and a niche definition fixed without using the target's acceptability.
The formal supplement states the corresponding observation model and its
identification assumptions. The agreement corpus probe reported there
illustrates the discipline: a rate per opportunity is evidence about choice,
and becomes evidence about licensing only after the alternatives and avoidance
option have been independently normed.

\subsection{The opening examples revisited}
\label{sec:worked-examples}

The profile now gives the opening asterisks different explanations. The results
are summarized in Table~\ref{tab:reader-profiles}; the labels describe the
pivotal condition rather than mutually exclusive sentence types.

\begin{table}[htbp]
\centering
\small
\begin{tabular}{@{}>{\raggedright\arraybackslash}p{0.27\textwidth}
  >{\raggedright\arraybackslash}p{0.24\textwidth}
  >{\raggedright\arraybackslash}p{0.38\textwidth}@{}}
\toprule
Case & Pivotal diagnosis & Expected profile \\
\midrule
\ungram{\mention{Can the have running}} & no covering analysis & confident rejection across construals \\
\mention{Colorless green ideas sleep furiously} & licensed but implausible content & grammatical status with content oddness \\
\ungram{\mention{I've finished it yesterday}} & incompatible temporal instructions & confident rejection resistant to exposure of the parts \\
\ungram{\mention{I have 25 years}} & concentrated non-licensing of the English age frame & recoverable meaning but confident, competitor-backed rejection \\
\marg{\mention{friend of whose}} & sparse opportunity and weak evidence & low confidence and comparatively strong framing or exposure effects \\
\ungram{\mention{Which did you buy car?}} & putative high-opportunity preemption & confident rejection only if the opportunity and competitor assumptions are independently met \\
\mention{pobedit'} first-person singular & winnerless cell & low rating with ineffability rather than confident correction \\
centre embedding & implementation overload & practice can improve the experience without changing licensing \\
\bottomrule
\end{tabular}
\caption{Reader-facing OVMG profiles. Each diagnosis makes a different
prediction beyond the judgment used to introduce it.}
\label{tab:reader-profiles}
\end{table}

\section{How licensing patterns change}
\label{sec:dynamics}

The state theory says what a present profile consists of. It doesn't explain how
the profile came to be. The dynamics module supplies candidate routes, with two
constraints on the explanation. First, learning about a community isn't the
same event as changing that community. Second, absence is informative only when
the missing candidate had a genuine opportunity to occur.

\subsection{Opportunity, candidates, and the outside option}
\label{sec:niches}
\label{sec:niche-sets}
\label{sec:usage}

A \term{niche} is a recurring communicative job, such as expressing a
recipient, locating an event in time, or marking an information source. Its
candidate set contains the realizations speakers could plausibly choose for
that job at that time. It also contains an outside option: paraphrase,
clarification, silence, or avoidance. The outside option matters because failure
to produce a target doesn't always mean that a competitor defeated it.

Production reflects both licensing and choice. A licensed form may be rare
because speakers prefer another licensed form. An unlicensed form may be absent
for the same surface reason. Conversely, a form may be scarcely attested simply
because the relevant niche is scarcely encountered. The analyst has to
establish the opportunity set, candidate preferences, and avoidance rate without
using the target's grammaticality judgment as the coding criterion.

The evidence stream distinguishes three events. A direct error provides evidence
about the target itself. An omission provides evidence only to the extent that
the target would plausibly have been chosen had it been available and that the
niche was correctly identified. Repair provides a further observation whose
meaning depends on its type and on the relation between the error and the public
update. These events enter one shared model but aren't interchangeable counts.

\subsection{Learning and population change}
\label{sec:update}

An individual updates an estimate when evidence arrives. The population changes
only when that evidence alters later inclusion, production, learning, or
transmission. This distinction is central. Updating an evidence-based estimate
can make beliefs more precise; it can't make speakers adopt or abandon a
construction. The formal dynamics add an explicit adoption-and-retention process
whose costs have to be motivated independently.

That process is intentionally conditional. Under some settings, weakly
supported constructions tend to disappear and strongly supported ones tend to
persist, producing separated population patterns. Under other settings, the
same learning rule yields a single interior pattern rather than separated
outer patterns.
The model doesn't derive categorical grammar from exposure alone.
It identifies the extra assumptions under which categorical-looking regimes can
emerge and makes those assumptions available for empirical challenge.

\subsection{What the simulations establish}
\label{sec:emergent-categoricality}

The simulations serve as existence and coherence checks. A deterministic
approximation locates candidate long-run regimes, and a stochastic agent
simulation checks whether finite populations with sampled evidence exhibit the
same qualitative behaviour. Together they show that separated licensing
patterns are possible when the adoption response, opportunity structure,
memory, and competition are jointly strong enough. They also show that the
separation disappears when those conditions are removed.

This conditional mechanism result isn't an empirical fit or a theorem that all
languages polarize. The model doesn't yet establish the exact long-run
behaviour of bounded-memory populations or how long they remain in one regime.
Nor do
the simulations currently model networked exposure, cohort replacement, or all
channels through a single shared underlying state. The supplement reports these
open proof and identification obligations alongside the successful checks.

\subsection{Repair as a candidate maintainer or controller}
\label{sec:repair-controller}

Repair can feed back into later learning, but different repairs do different
work. Some secure understanding without endorsing the local form. Some replace
an exponent but leave the intended operator value intact. Some correct an
operator mis-setting that would otherwise change who is committed to what,
when, or with which discourse status. Overt correction can also reflect
prescriptive or identity policing rather than a threat to the update.

The operator hypothesis predicts a selective response profile. Once closure,
opportunity, competitor strength, and ordinary comprehensibility are controlled,
mis-setting an update-configuring contrast should disproportionately elicit
update-oriented repair. A categorically unlicensed but non-operator construction
may instead be corrected as a conventional form choice. The prediction concerns
the repair target. Non-operators can still be categorically policed.

Calling repair a \emph{controller} requires more than observing that correction
and stability coexist. An intervention would have to show that the system
registers the relevant deviation, changes lower-level behaviour in response,
and thereby restores the threatened population relation. Evidence that removing
repair weakens transmission would establish maintenance; evidence of the full
deviation-sensitive loop would establish control.

\subsection{Four predicted trajectories}

The dynamics yield four qualitatively different trajectories that would be
collapsed by a single frequency measure.

\noindent\textsc{Obligatorification.}
\label{sec:derived-obligatoriness}
An overt contrast becomes obligatory when unmarked realizations repeatedly lose
in well-attested niches and that pattern persists long enough to become a
community convention. English progressive marking and Turkish evidential
marking are proposed cases. The claim is historical and niche-specific: the
model can't infer obligatoriness from the analyst's present uncertainty or
from the absence of data. Once established, obligatoriness also doesn't vanish
with the first contrary token; recovery has to persist before the convention
changes.

\noindent\textsc{Winnerless cells.}
\label{sec:winnerless-cells}
A dense niche can lack a conventional winner. Speakers may paraphrase or avoid
the job. Because avoidance isn't readily attributable to any one candidate, it
doesn't by itself lower support for a particular candidate; it mainly starves
all of them of evidence. Low licensing with low confidence also requires
independently motivated weak starting support or diffuse negative
evidence. The resulting profile is ineffability rather than a competitor-backed
ban. Russian defective paradigms supply the clean model; English
\ungram{\mention{amn't}} is likely mixed because a competitor and avoidance both
matter.

\noindent\textsc{Destabilization.}
\label{sec:destabilization}
A moribund contrast can lose evidential concentration before its average status
changes much. If a shared evidence stream disappears, individual histories can
drift back toward starting expectations. Between-speaker divergence follows
early only when a further source of heterogeneity (different networks, cohorts,
starting expectations, or exposure histories) is present. Concentration loss is the direct
prediction; population dispersion is conditional.

\noindent\textsc{Actuation and apparent categorical exclusions.}
\label{sec:actuation}
\label{sec:categorical-without-bans}
Small changes in opportunity, competitor strength, analogy, repair, or social
ascription can move a construction across a boundary between long-run regimes.
That supplies candidate actuation routes without pretending to predict the
historical trigger from the state theory alone. At the other extreme, persistent
preemption in a dense niche can generate a highly concentrated rejection even
when a compatible analysis remains available. Apparent categorical exclusions
don't always require an additional structural ban; they require
independent evidence that the niche was rich, the target was a live candidate,
and a competitor repeatedly won.

The dynamics section's empirical burden is now visible without the equations.
The account succeeds only if independently measured opportunities, alternatives,
repair types, evidence histories, and adoption costs predict the distinct
trajectories. The formal supplement supplies the quantitative claims and
reproducible simulation contract; the main paper keeps their linguistic meaning
and their failure conditions in view.
